% =====================================================================
% Epilogue --- Duality
% Scenes 25--27: Compression-Prediction Twin, KL as protagonist, Funes returns.
% =====================================================================

\renewcommand{\AINITsection}{Epilogue \textemdash{} Duality}

\sectiondividerwithportrait%
  {hutter.jpg}%
  {``Compression is comprehension.''}%
  {Marcus Hutter, on the spirit of his prize}%
  {Epilogue \textemdash{} Duality}

% ---------------------------------------------------------------------
% Scene 25a --- Compression-Prediction Twin (framing slide)
% ---------------------------------------------------------------------
\begin{frame}{A perfect predictor is a perfect compressor}
  Two halves of one subject.

  \begin{itemize}
    \item A \AINITemph{forecaster} assigns probabilities to the next symbol;
          her loss is the average $-\log \hat p$ she pays per draw.
    \item An \AINITemph{encoder} assigns short codewords to likely symbols;
          his cost is the average bits per draw.
    \item Shannon (1948) proved these two numbers are the same number.
          Both bottom out at $H(p)$.
  \end{itemize}

  \vspace{0.5em}
  Hence the slogan, due to Hutter\textsuperscript{1}:
  \AINITemph{compression is comprehension.}
  And the empirical confirmation, due to Del\'etang et al.\textsuperscript{2}:
  a large language model, trained only on text, compresses ImageNet
  pixels better than PNG.

  \vfill
  {\footnotesize
   \textsuperscript{1}\,Hutter, M. (2005). \emph{Universal Artificial Intelligence};
   the Hutter Prize for lossless text compression.\\
   \textsuperscript{2}\,Del\'etang, G.\ et al.\ (2023).
   \emph{Language Modeling Is Compression}.}
\end{frame}

% ---------------------------------------------------------------------
% Scene 25b --- viz handoff
% ---------------------------------------------------------------------
\vizhandoff{compression-prediction-twin.html}

% ---------------------------------------------------------------------
% Scene 26 --- Cross-entropy = floor + inefficiency
% ---------------------------------------------------------------------
\equationslide%
  {Floor + inefficiency}%
  {%
    \underbrace{-\mathbb{E}_p[\log \hat p]}_{\text{cross-entropy}}
    \;=\;
    \underbrace{H(p)}_{\text{Shannon's floor}}
    \;+\;
    \underbrace{{\color{accentOxblood}\mathrm{KL}(p \,\|\, \hat p)}}_{\text{your inefficiency}}
  }%
  {Recsys: KL is the gap between your factorization and the truth.
   Bandits: KL is the cost of mistaking one urn for another.
   Compression: KL is the bits you waste because your predictor is imperfect.
   One quantity, three roles.}

% ---------------------------------------------------------------------
% Scene 27 --- Funes returns (closing slide)
% ---------------------------------------------------------------------
\closingslide%
  {Funes remembered everything and could think nothing.}%
  {Machine learning forgets almost everything --- and that is why it understands.}
